\chapter{其他职业病}
\setcounter{chapter}{8}
\chapterintro{
\textbf{【教学大纲要求】}

本章为\imp{自学}内容。包括职业性皮肤病、职业性眼病、职业性耳鼻喉及口腔疾病等。
}

\section{职业性皮肤病}

\begin{defbox}
\textbf{职业性皮肤病}：在职业活动中接触化学、物理、生物等职业性有害因素引起的皮肤及其附属器的疾病。是职业病中最常见的一大类。

\textbf{常见类型：}
\begin{enumerate}[leftmargin=*]
    \item \imp{职业性皮炎}：刺激性接触性皮炎（最常见）和变应性接触性皮炎
    \item \imp{职业性皮肤色素变化}：色素沉着（煤焦油、沥青）和色素减退
    \item \imp{职业性痤疮}：油痤疮（矿物油）和氯痤疮（卤代芳烃如二噁英）
    \item \imp{职业性皮肤溃疡}：铬疮（六价铬化合物引起的深在性溃疡）
    \item \imp{职业性皮肤肿瘤}：煤焦油和沥青→皮肤癌
    \item \imp{职业性感染性皮肤病}：炭疽等
\end{enumerate}
\end{defbox}

\section{职业性眼病}

\begin{itemize}[leftmargin=*]
    \item \imp{化学性眼灼伤}：酸碱溅入眼内
    \item \imp{电光性眼炎}：紫外线照射→急性角膜结膜炎（潜伏期4--12h）
    \item \imp{职业性白内障}：红外线（玻璃工人白内障）、TNT、微波等
\end{itemize}

\section{职业性耳鼻喉及口腔疾病}

\begin{itemize}[leftmargin=*]
    \item \imp{噪声聋}：感音性耳聋，高频听力下降为主
    \item \imp{铬鼻病}：六价铬引起鼻中隔穿孔
    \item \imp{牙酸蚀症}：长期接触酸雾导致牙齿硬组织脱钙
\end{itemize}

\newpage

% =====================================================================
%  CHAPTER 9: 职业伤害（职业安全）
% =====================================================================
\newpage